\documentclass[10pt,a4paper]{amsart}
\usepackage{amsmath,amssymb,amsthm,graphicx}
% Retained minimal preamble for the Stacks Project chapter flat.tex (More on Flatness).
% Verbatim excerpt of the chapter preamble: Xy-pic setup (native preamble lines 14-18) and the native
% theorem declarations with the mathematical macros (native preamble lines 154-217). Copied mechanically;
% nothing added, removed or redefined. The wrapper class and the amsmath/amssymb/amsthm/graphicx base
% packages are supplied by the assembler, so the class/externaldocument part of the native preamble is omitted.
% For commutative diagrams we use Xy-pic
\usepackage[all]{xy}

% We use 2cell for 2-commutative diagrams.
\xyoption{2cell}
% Theorem environments.
%
\theoremstyle{plain}
\newtheorem{theorem}[subsection]{Theorem}
\newtheorem{proposition}[subsection]{Proposition}
\newtheorem{lemma}[subsection]{Lemma}

\theoremstyle{definition}
\newtheorem{definition}[subsection]{Definition}
\newtheorem{example}[subsection]{Example}
\newtheorem{exercise}[subsection]{Exercise}
\newtheorem{situation}[subsection]{Situation}

\theoremstyle{remark}
\newtheorem{remark}[subsection]{Remark}
\newtheorem{remarks}[subsection]{Remarks}

\numberwithin{equation}{subsection}

% Macros
%
\def\lim{\mathop{\mathrm{lim}}\nolimits}
\def\colim{\mathop{\mathrm{colim}}\nolimits}
\def\Spec{\mathop{\mathrm{Spec}}}
\def\Hom{\mathop{\mathrm{Hom}}\nolimits}
\def\Ext{\mathop{\mathrm{Ext}}\nolimits}
\def\SheafHom{\mathop{\mathcal{H}\!\mathit{om}}\nolimits}
\def\SheafExt{\mathop{\mathcal{E}\!\mathit{xt}}\nolimits}
\def\Sch{\mathit{Sch}}
\def\Mor{\mathop{\mathrm{Mor}}\nolimits}
\def\Ob{\mathop{\mathrm{Ob}}\nolimits}
\def\Sh{\mathop{\mathit{Sh}}\nolimits}
\def\NL{\mathop{N\!L}\nolimits}
\def\CH{\mathop{\mathrm{CH}}\nolimits}
\def\proetale{{pro\text{-}\acute{e}tale}}
\def\etale{{\acute{e}tale}}
\def\QCoh{\mathit{QCoh}}
\def\Ker{\mathop{\mathrm{Ker}}}
\def\Im{\mathop{\mathrm{Im}}}
\def\Coker{\mathop{\mathrm{Coker}}}
\def\Coim{\mathop{\mathrm{Coim}}}

% Boxtimes
%
\DeclareMathSymbol{\boxtimes}{\mathbin}{AMSa}{"02}

%
% Macros for moduli stacks/spaces
%
\def\QCohstack{\mathcal{QC}\!\mathit{oh}}
\def\Cohstack{\mathcal{C}\!\mathit{oh}}
\def\Spacesstack{\mathcal{S}\!\mathit{paces}}
\def\Quotfunctor{\mathrm{Quot}}
\def\Hilbfunctor{\mathrm{Hilb}}
\def\Curvesstack{\mathcal{C}\!\mathit{urves}}
\def\Polarizedstack{\mathcal{P}\!\mathit{olarized}}
\def\Complexesstack{\mathcal{C}\!\mathit{omplexes}}
% \Pic is the operator that assigns to X its picard group, usage \Pic(X)
% \Picardstack_{X/B} denotes the Picard stack of X over B
% \Picardfunctor_{X/B} denotes the Picard functor of X over B
\def\Pic{\mathop{\mathrm{Pic}}\nolimits}
\def\Picardstack{\mathcal{P}\!\mathit{ic}}
\def\Picardfunctor{\mathrm{Pic}}
\def\Deformationcategory{\mathcal{D}\!\mathit{ef}}

% --- appended 2026-10-12 (Chief): the project own definitions of the reference/slogan/history
% display environments, copied mechanically from the shared Stacks preamble lines 8-12
% (\usepackage{verbatim} plus the three \newenvironment lines, which the chapter preamble keeps
% verbatim from the project preamble). Nothing added, removed or redefined: these environments
% carry no mathematical content and the project itself typesets them as comments. Without these
% three lines a counted statement whose own text uses them fails with Environment slogan undefined.
\usepackage{verbatim}
\newenvironment{reference}{\comment}{\endcomment}
\newenvironment{slogan}{\comment}{\endcomment}
\newenvironment{history}{\comment}{\endcomment}
\usepackage{cleveref}
\crefname{originaldefinition}{original definition}{original definitions}
\Crefname{originaldefinition}{Original definition}{Original definitions}
\makeatletter
\newlabel{algebra-definition-henselization}{{04GQ}{source}{}{sourcepointer-1}{}}
\newlabel{algebra-definition-henselization@cref}{{[originaldefinition][0][]04GQ}{[1][0][]source}{}{}{}}
\makeatother
\crefname{originallemma}{original lemma}{original lemmas}
\Crefname{originallemma}{Original lemma}{Original lemmas}
\makeatletter
\newlabel{algebra-lemma-colimit-rings-flat}{{05UU}{source}{}{sourcepointer-2}{}}
\newlabel{algebra-lemma-colimit-rings-flat@cref}{{[originallemma][0][]05UU}{[1][0][]source}{}{}{}}
\makeatother
\crefname{originallemma}{original lemma}{original lemmas}
\Crefname{originallemma}{Original lemma}{Original lemmas}
\makeatletter
\newlabel{algebra-lemma-composition-flat}{{00HC}{source}{}{sourcepointer-3}{}}
\newlabel{algebra-lemma-composition-flat@cref}{{[originallemma][0][]00HC}{[1][0][]source}{}{}{}}
\makeatother
\crefname{originallemma}{original lemma}{original lemmas}
\Crefname{originallemma}{Original lemma}{Original lemmas}
\makeatletter
\newlabel{algebra-lemma-flatness-descends}{{00HJ}{source}{}{sourcepointer-4}{}}
\newlabel{algebra-lemma-flatness-descends@cref}{{[originallemma][0][]00HJ}{[1][0][]source}{}{}{}}
\makeatother
\crefname{originallemma}{original lemma}{original lemmas}
\Crefname{originallemma}{Original lemma}{Original lemmas}
\makeatletter
\newlabel{algebra-lemma-henselization-different}{{04GV}{source}{}{sourcepointer-5}{}}
\newlabel{algebra-lemma-henselization-different@cref}{{[originallemma][0][]04GV}{[1][0][]source}{}{}{}}
\makeatother
\crefname{originallemma}{original lemma}{original lemmas}
\Crefname{originallemma}{Original lemma}{Original lemmas}
\makeatletter
\newlabel{algebra-lemma-strict-henselization-different}{{04GW}{source}{}{sourcepointer-6}{}}
\newlabel{algebra-lemma-strict-henselization-different@cref}{{[originallemma][0][]04GW}{[1][0][]source}{}{}{}}
\makeatother
\crefname{originallemma}{original lemma}{original lemmas}
\Crefname{originallemma}{Original lemma}{Original lemmas}
\makeatletter
\newlabel{lemma-etale-flat-up-down-local-ring}{{05BA}{source}{}{sourcepointer-7}{}}
\newlabel{lemma-etale-flat-up-down-local-ring@cref}{{[originallemma][0][]05BA}{[1][0][]source}{}{}{}}
\makeatother
\crefname{originallemma}{original lemma}{original lemmas}
\Crefname{originallemma}{Original lemma}{Original lemmas}
\makeatletter
\newlabel{more-algebra-lemma-dumb-properties-henselization}{{07QM}{source}{}{sourcepointer-8}{}}
\newlabel{more-algebra-lemma-dumb-properties-henselization@cref}{{[originallemma][0][]07QM}{[1][0][]source}{}{}{}}
\makeatother
\title{Flatness is the same over a local ring, over its henselization and over its strict henselization}
\author{The Stacks Project Authors}
\date{stacks-project:flat.tex:local-snapshot-20261011; GNU FDL 1.2 or later}
\begin{document}
\maketitle
\noindent\textit{Editorial note.} All delivered mathematical text is the author's own, copied from the pinned Stacks Project chapter file flat.tex; no statement or proof step was rewritten, reordered or removed.\par\medskip
\noindent\textit{Editorial note.} Preamble provenance: the chapter's own retained minimal preamble (sources/stacks-flat/retained-preamble.tex, lines 1-85) is a verbatim excerpt of the native Xy-pic setup and of the native theorem declarations and mathematical macros, followed by the project's own reference/slogan/history display-environment definitions copied mechanically from the shared preamble; the wrapper class and the amsmath/amssymb/amsthm/graphicx base packages are supplied by the assembler, and extra\_packages is left empty because everything needed is already in the retained preamble. No macro was invented or redefined.\par\medskip
\noindent\textit{Editorial note.} Same-chapter pointers are located at their \textbackslash{}label lines inside the pinned chapter file; cross-chapter pointers are located in the project's own tag list staged in the snapshot (sources/stacks-chow/tags.txt), at the verbatim lines carrying their canonical tags.\par\medskip
\noindent\textit{Allowed standard prerequisites (editorial).}\begin{itemize}
\item Henselization and strict henselization and their faithful flatness over the local ring. Algebra, label algebra-definition-henselization (statement, proof line 187); More on Algebra, label more-algebra-lemma-dumb-properties-henselization (proof line 210). Fixes the rings and gives that the local ring maps from the local ring to its henselization and to its strict henselization are faithfully flat, which yields the implications downwards by composition of flat maps.
\item Composition and descent of flatness along faithfully flat ring maps. Algebra, labels algebra-lemma-composition-flat (proof line 216) and algebra-lemma-flatness-descends (proof line 219). The first gives the implications (3) implies (2) implies (1) and (5) implies (4), and the second gives the equivalences relating the statements about a module over the local ring or over the henselization to the statement over the strict henselization.
\item Flatness up and down along an etale neighbourhood, the colimit descriptions of the henselizations, and flatness of a colimit. Same chapter flat.tex, label lemma-etale-flat-up-down-local-ring (proof line 227); Algebra, labels algebra-lemma-henselization-different and algebra-lemma-strict-henselization-different (proof lines 232-233) and algebra-lemma-colimit-rings-flat (proof line 236). Together they show that flatness over the local ring persists to every etale neighbourhood and hence, because the henselizations are directed colimits of those local rings, to the henselization and the strict henselization, which proves the remaining implication.
\end{itemize}\medskip
\ifdefined\hypertarget\hypertarget{sourcepointer-1}{}\fi
\noindent\textit{Source reference.} Original definition 04GQ: sources/stacks-chow/tags.txt, source lines 5841--5841. Navigation pointer for \textbackslash{}ref\{algebra-definition-henselization\} used inside the counted span (line 187); the argument that consumes it is retained in the span.\par\medskip
\ifdefined\hypertarget\hypertarget{sourcepointer-2}{}\fi
\noindent\textit{Source reference.} Original lemma 05UU: sources/stacks-chow/tags.txt, source lines 7528--7528. Navigation pointer for \textbackslash{}ref\{algebra-lemma-colimit-rings-flat\} used inside the counted span (line 236); the argument that consumes it is retained in the span.\par\medskip
\ifdefined\hypertarget\hypertarget{sourcepointer-3}{}\fi
\noindent\textit{Source reference.} Original lemma 00HC: sources/stacks-chow/tags.txt, source lines 638--638. Navigation pointer for \textbackslash{}ref\{algebra-lemma-composition-flat\} used inside the counted span (line 216); the argument that consumes it is retained in the span.\par\medskip
\ifdefined\hypertarget\hypertarget{sourcepointer-4}{}\fi
\noindent\textit{Source reference.} Original lemma 00HJ: sources/stacks-chow/tags.txt, source lines 645--645. Navigation pointer for \textbackslash{}ref\{algebra-lemma-flatness-descends\} used inside the counted span (line 219); the argument that consumes it is retained in the span.\par\medskip
\ifdefined\hypertarget\hypertarget{sourcepointer-5}{}\fi
\noindent\textit{Source reference.} Original lemma 04GV: sources/stacks-chow/tags.txt, source lines 5846--5846. Navigation pointer for \textbackslash{}ref\{algebra-lemma-henselization-different\} used inside the counted span (line 232); the argument that consumes it is retained in the span.\par\medskip
\ifdefined\hypertarget\hypertarget{sourcepointer-6}{}\fi
\noindent\textit{Source reference.} Original lemma 04GW: sources/stacks-chow/tags.txt, source lines 5847--5847. Navigation pointer for \textbackslash{}ref\{algebra-lemma-strict-henselization-different\} used inside the counted span (line 233); the argument that consumes it is retained in the span.\par\medskip
\ifdefined\hypertarget\hypertarget{sourcepointer-7}{}\fi
\noindent\textit{Source reference.} Original lemma 05BA: sources/stacks-flat/pinned-original-source.tex, source lines 160--160. Navigation pointer for \textbackslash{}ref\{lemma-etale-flat-up-down-local-ring\} used inside the counted span (line 227); the argument that consumes it is retained in the span.\par\medskip
\ifdefined\hypertarget\hypertarget{sourcepointer-8}{}\fi
\noindent\textit{Source reference.} Original lemma 07QM: sources/stacks-chow/tags.txt, source lines 9837--9837. Navigation pointer for \textbackslash{}ref\{more-algebra-lemma-dumb-properties-henselization\} used inside the counted span (line 210); the argument that consumes it is retained in the span.\par\medskip
\begin{lemma}
\label{lemma-flat-up-down-henselization}
Let $S$ be a scheme and $s \in S$ a point. Denote $\mathcal{O}_{S, s}^h$
(resp.\ $\mathcal{O}_{S, s}^{sh}$) the henselization (resp.\ strict
henselization), see
Algebra, Definition \ref{algebra-definition-henselization}.
Let $M^{sh}$ be a $\mathcal{O}_{S, s}^{sh}$-module.
The following are equivalent
\begin{enumerate}
\item $M^{sh}$ is flat over $\mathcal{O}_{S, s}$,
\item $M^{sh}$ is flat over $\mathcal{O}_{S, s}^h$, and
\item $M^{sh}$ is flat over $\mathcal{O}_{S, s}^{sh}$.
\end{enumerate}
If $M^{sh} = M^h \otimes_{\mathcal{O}_{S, s}^h} \mathcal{O}_{S, s}^{sh}$
this is also equivalent to
\begin{enumerate}
\item[(4)] $M^h$ is flat over $\mathcal{O}_{S, s}$, and
\item[(5)] $M^h$ is flat over $\mathcal{O}_{S, s}^h$.
\end{enumerate}
If $M^h = M \otimes_{\mathcal{O}_{S, s}} \mathcal{O}_{S, s}^h$
this is also equivalent to
\begin{enumerate}
\item[(6)] $M$ is flat over $\mathcal{O}_{S, s}$.
\end{enumerate}
\end{lemma}

\begin{proof}
By More on Algebra, Lemma
\ref{more-algebra-lemma-dumb-properties-henselization}
the local ring maps
$\mathcal{O}_{S, s} \to \mathcal{O}_{S, s}^h \to \mathcal{O}_{S, s}^{sh}$
are faithfully flat.
Hence (3) $\Rightarrow$ (2) $\Rightarrow$ (1) and
(5) $\Rightarrow$ (4) follow from
Algebra, Lemma \ref{algebra-lemma-composition-flat}.
By faithful flatness the equivalences (6) $\Leftrightarrow$ (5) and
(5) $\Leftrightarrow$ (3) follow from
Algebra, Lemma \ref{algebra-lemma-flatness-descends}.
Thus it suffices to show that
(1) $\Rightarrow$ (2) $\Rightarrow$ (3) and
(4) $\Rightarrow$ (5).
To prove these we may assume $S$ is an affine scheme.

\medskip\noindent
Assume (1). By
Lemma \ref{lemma-etale-flat-up-down-local-ring}
we see that $M^{sh}$ is flat over $\mathcal{O}_{T, t}$ for
any \'etale neighbourhood $(T, t) \to (S, s)$. Since $\mathcal{O}_{S, s}^h$
and $\mathcal{O}_{S, s}^{sh}$ are directed colimits of local rings
of the form $\mathcal{O}_{T, t}$ (see
Algebra, Lemmas \ref{algebra-lemma-henselization-different}
and \ref{algebra-lemma-strict-henselization-different})
we conclude that $M^{sh}$ is flat over $\mathcal{O}_{S, s}^h$
and $\mathcal{O}_{S, s}^{sh}$ by
Algebra, Lemma \ref{algebra-lemma-colimit-rings-flat}.
Thus (1) implies (2) and (3). Of course this implies also
(2) $\Rightarrow$ (3) by replacing $\mathcal{O}_{S, s}$ by
$\mathcal{O}_{S, s}^h$. The same argument applies to prove
(4) $\Rightarrow$ (5).
\end{proof}

\end{document}
